% ============================================================
%  INFORME TÉCNICO — Práctica Experimental Unidad IV
%  GA — Aplicaciones Web  |  Grupo: Pallo · Beltrán · Cruz
% ============================================================
\documentclass[12pt,a4paper]{article}

% ── Codificación y fuentes ──────────────────────────────────
\usepackage[utf8]{inputenc}
\usepackage[T1]{fontenc}
\usepackage{lmodern}
\usepackage{microtype}

% ── Idioma ─────────────────────────────────────────────────
\usepackage[spanish,english]{babel}

% ── Geometría ──────────────────────────────────────────────
\usepackage[
  top=2.5cm, bottom=2.5cm,
  left=3cm, right=2.5cm
]{geometry}

% ── Interlineado ───────────────────────────────────────────
\usepackage{setspace}
\setstretch{1.5}

% ── Encabezados y pies ─────────────────────────────────────
\usepackage{fancyhdr}
\pagestyle{fancy}
\fancyhf{}
\fancyhead[L]{\small Práctica Experimental --- Unidad IV}
\fancyhead[R]{\small Pallo $\cdot$ Beltrán $\cdot$ Cruz}
\fancyfoot[C]{\thepage}
\renewcommand{\headrulewidth}{0.4pt}

% ── Tablas ─────────────────────────────────────────────────
\usepackage{booktabs}
\usepackage{tabularx}
\usepackage{array}
\newcolumntype{L}[1]{>{\raggedright\arraybackslash}p{#1}}

% ── Colores y cajas ────────────────────────────────────────
\usepackage[dvipsnames]{xcolor}
\usepackage{amssymb}
\usepackage{tcolorbox}
\tcbuselibrary{skins,breakable}

% ── Código fuente ──────────────────────────────────────────
\usepackage{caption}
\usepackage{listings}
\lstdefinestyle{json}{
  language=,
  basicstyle=\ttfamily\small,
  breaklines=true,
  frame=single,
  backgroundcolor=\color{gray!8},
  rulecolor=\color{gray!40},
  numbers=none,
  showstringspaces=false
}
\lstdefinestyle{xml}{
  language=XML,
  basicstyle=\ttfamily\small,
  breaklines=true,
  frame=single,
  backgroundcolor=\color{orange!5},
  rulecolor=\color{orange!40},
  numbers=none,
  showstringspaces=false
}
\lstdefinestyle{yaml}{
  language=,
  basicstyle=\ttfamily\small,
  breaklines=true,
  frame=single,
  backgroundcolor=\color{blue!4},
  rulecolor=\color{blue!20},
  numbers=none,
  showstringspaces=false
}

% ── Gráficos e imágenes ────────────────────────────────────
\usepackage{graphicx}


% ── Hipervínculos ──────────────────────────────────────────
\usepackage[
  colorlinks=true,
  linkcolor=NavyBlue,
  citecolor=ForestGreen,
  urlcolor=BrickRed
]{hyperref}

% ── Citas APA (biblatex + biber) ───────────────────────────
\usepackage[
  style=apa,
  sortcites=true,
  sorting=nyt,
  backend=biber
]{biblatex}
\addbibresource{referencias.bib}

% ── Formato de secciones ───────────────────────────────────
\usepackage{titlesec}
\titleformat{\section}{\large\bfseries\color{NavyBlue}}{\thesection}{1em}{}
\titleformat{\subsection}{\normalsize\bfseries\color{NavyBlue!80}}{\thesubsection}{1em}{}

% ── Metadatos PDF ──────────────────────────────────────────
\hypersetup{
  pdftitle={Práctica Experimental Unidad IV},
  pdfauthor={Pallo, Beltrán, Cruz},
  pdfsubject={Aplicaciones Web --- Unidad IV: Temas 13-16},
  pdfkeywords={REST, SOAP, MVC, API, Swagger, Jamstack, PWA}
}

% ============================================================
\begin{document}
\selectlanguage{spanish}

% ────────────────────────────────────────────────────────────
%  PORTADA
% ────────────────────────────────────────────────────────────
\begin{titlepage}
  \begin{center}
    \vspace*{1cm}
    {\LARGE \textbf{UNIVERSIDAD TÉCNICA}}\\[0.3cm]
    {\large Facultad de Ciencias de la Computación}\\[0.2cm]
    {\large Carrera de Ingeniería en Tecnologías de la Información}\\[2cm]

    \rule{\linewidth}{0.5mm}\\[0.4cm]
    {\huge \textbf{Práctica Experimental --- Unidad IV}}\\[0.3cm]
    {\Large Temas 13--16: Arquitecturas Web, APIs REST,\\[0.1cm]
     SOAP y Tendencias Emergentes}\\[0.3cm]
    \rule{\linewidth}{0.5mm}\\[1.5cm]

    {\large \textbf{Materia:} Aplicaciones Web}\\[0.3cm]
    {\large \textbf{Grupo:} PFC --- Pallo $\cdot$ Beltrán $\cdot$ Cruz}\\[0.3cm]
    {\large \textbf{PFC:} SGED --- Sistema de Gestión Estudiantil Deportiva}\\[2cm]

    \vfill
    {\large \today}
  \end{center}
\end{titlepage}

\newpage

% ────────────────────────────────────────────────────────────
%  RESUMEN BILINGÜE  [Cruz]
% ────────────────────────────────────────────────────────────
\section*{Resumen}
\addcontentsline{toc}{section}{Resumen}

\noindent
El presente informe documenta la Práctica Experimental correspondiente a la Unidad IV de la
asignatura Aplicaciones Web, abarcando los Temas 13 a 16.
El equipo diseñó, implementó y demostró el \textbf{Sistema de Gestión Estudiantil Deportiva (SGED)},
una aplicación web completa basada en el patrón Modelo--Vista--Controlador (MVC) con backend en
\textit{Spring Boot 3.2} (Java~21) y frontend en \textit{Angular 21}.
Se expone la API REST del sistema con documentación interactiva mediante Swagger UI accesible en
\texttt{/api/docs}, se ejecutan en vivo al menos tres endpoints representativos y se elabora una
comparativa rigurosa entre SOAP y REST con nueve criterios de análisis.
Finalmente, se reflexiona críticamente sobre las tendencias emergentes: Jamstack,
Progressive Web Apps (PWA) e Inteligencia Artificial generativa aplicada al desarrollo web.

\bigskip

\begin{otherlanguage}{english}
\noindent\textbf{Abstract.}
This report documents the Experimental Practice for Unit IV of the Web Applications course,
covering Topics 13--16.
The team designed, implemented, and demonstrated the \textbf{Sports Student Management System (SGED)},
a full-stack web application built on the Model--View--Controller (MVC) pattern using a
\textit{Spring Boot 3.2} (Java~21) backend and an \textit{Angular 21} frontend.
The system's REST API is presented with interactive documentation via Swagger UI at \texttt{/api/docs};
at least three representative endpoints are executed live.
A rigorous nine-criterion comparison between SOAP and REST is provided, and the report closes with a
critical reflection on emerging trends: Jamstack, Progressive Web Apps (PWA), and generative
Artificial Intelligence in web development.
\end{otherlanguage}

\bigskip
\noindent\textbf{Palabras clave:} REST, SOAP, MVC, API, OpenAPI, Swagger, Jamstack, PWA,
IA generativa, Spring Boot.\\
\textbf{Keywords:} REST, SOAP, MVC, API, OpenAPI, Swagger, Jamstack, PWA, generative AI,
Spring Boot.

\newpage
\tableofcontents
\newpage

% ============================================================
%  1. INTRODUCCIÓN  [Cruz]
% ============================================================
\section{Introducción}

El ecosistema del desarrollo web ha evolucionado aceleradamente durante la última década.
La adopción masiva de arquitecturas orientadas a servicios, el auge de las API REST y la irrupción
de paradigmas como Jamstack y la IA generativa obligan al profesional contemporáneo a comprender
tanto los fundamentos teóricos como sus implicaciones prácticas.

Esta práctica tiene como propósito consolidar los conocimientos de la Unidad IV a través de la
implementación, documentación y análisis crítico de un sistema real: el SGED
(\textit{Sports Student Management System}), desarrollado como Proyecto de Fin de Carrera (PFC)
del equipo.

El informe se organiza como sigue: la Sección~\ref{sec:biblio} presenta la revisión bibliográfica
con al menos cinco fuentes académicas en norma APA; la Sección~\ref{sec:mvc} describe la
aplicación MVC completa; la Sección~\ref{sec:api} documenta la API REST con Swagger UI;
la Sección~\ref{sec:soap} elabora la comparativa SOAP vs.\ REST; la Sección~\ref{sec:tendencias}
reflexiona sobre las tendencias emergentes; y las Secciones~\ref{sec:conclusiones}
y~\ref{sec:futuro} presentan las conclusiones y el trabajo futuro, respectivamente.

% ============================================================
%  2. REVISIÓN BIBLIOGRÁFICA  [Beltrán]  — Directriz 1
% ============================================================
\section{Revisión Bibliográfica}
\label{sec:biblio}

Como parte de la Directriz~1, esta sección revisa seis fuentes académicas y técnicas primarias
--- superando el mínimo de cinco exigido --- y las contrasta directamente contra el código y el
comportamiento observado del SGED durante la verificación en vivo documentada en la
Sección~\ref{sec:mvc}. La disertación de \textcite{fielding2000} y la documentación oficial de
\textcite{springdoc2024} se incluyen por ser, respectivamente, la fuente seminal de REST y la
librería que el propio backend usa para exponer su documentación.

\subsection{Fielding (2000) --- REST como estilo arquitectónico}

\textcite{fielding2000} define REST no como un protocolo sino como un conjunto de seis
restricciones arquitectónicas --- cliente--servidor, \textit{stateless}, cacheable, interfaz
uniforme, sistema por capas y código bajo demanda (opcional) --- cuya combinación produce
sistemas escalables y evolucionables sin acoplar cliente y servidor a una implementación
concreta \parencite{fielding2000}. Verificando el SGED contra estas seis restricciones:
la autenticación es efectivamente \textit{stateless} (el JWT viaja en una \textit{cookie}
\texttt{HttpOnly}, no en sesión de servidor); la interfaz es mayormente uniforme
(\texttt{/api/categorias}, \texttt{/api/entrenadores}, sustantivos en plural); pero la
restricción de \textit{cacheable} no se cumple de forma consistente --- no se observaron
cabeceras \texttt{Cache-Control} en las respuestas inspeccionadas por Swagger durante la
verificación en vivo --- lo cual es una brecha real entre la teoría de Fielding y la
implementación, no solo una observación de manual.

\subsection{Richardson y Ruby (2007) --- REST Práctico}

\textcite{richardson2007} amplían a Fielding con una guía operativa de diseño de APIs
RESTful, introduciendo la \textit{Resource-Oriented Architecture} (ROA) y un modelo de
madurez por niveles (Richardson Maturity Model). Bajo ese modelo, el SGED alcanza el
Nivel~2 (recursos + verbos HTTP correctos) pero no el Nivel~3 (HATEOAS): las respuestas de
\texttt{GET /api/categorias/\{id\}} devuelven el recurso sin enlaces a acciones relacionadas
(por ejemplo, a las sesiones de esa categoría), algo que sí exigiría un API verdaderamente
madura según \textcite{richardson2007}.

\subsection{SpringDoc OpenAPI 2 --- Documentación oficial}

\textcite{springdoc2024} documenta la integración de SpringDoc OpenAPI~2 con Spring Boot~3,
la librería que el SGED usa para exponer su especificación OpenAPI en
\texttt{/api/docs.json} y la interfaz interactiva Swagger UI en \texttt{/api/docs}. Esta
configuración no es la ruta por defecto de la librería (que sería
\texttt{/swagger-ui/index.html}), sino un ajuste explícito del equipo:

\begin{lstlisting}[style=yaml, caption={Configuración SpringDoc en \texttt{application.yml}}]
springdoc:
  api-docs:
    path: /api/docs.json
  swagger-ui:
    path: /api/docs
    operationsSorter: method
\end{lstlisting}

Verificado en vivo: acceder a la ruta por defecto de SpringDoc produce el error
\textit{``Failed to load remote configuration''}, porque la especificación OpenAPI
efectivamente vive en \texttt{/api/docs.json}, no en la ruta que SpringDoc asume por
convención. Es un detalle de configuración correcto pero no evidente para quien no conoce
de antemano el ajuste --- vale la pena documentarlo explícitamente en el \texttt{README} del
proyecto, algo que actualmente no ocurre.

\subsection{W3C (2007) --- Especificación SOAP 1.2}

\textcite{w3c2007} normaliza el protocolo SOAP~1.2: mensajería basada en XML con un
\textit{Envelope} obligatorio (\textit{Header} + \textit{Body}), independiente del
transporte subyacente, y descrita formalmente mediante WSDL. Esta fuente sustenta la
comparativa SOAP vs.\ REST desarrollada por Pallo en la Sección~\ref{sec:soap}; se cita aquí
porque el contraste entre el tipado fuerte y \textit{machine-readable} de WSDL frente al
contrato implícito (y, como se documenta en la Sección~\ref{sec:mvc}, a veces inconsistente)
de la API REST del SGED es un argumento recurrente en el análisis crítico de esta unidad.

\subsection{Masse (2011) --- REST API Design Rulebook}

\textcite{masse2011} sistematiza reglas prescriptivas de diseño REST: nomenclatura de URIs
en minúsculas con sustantivos en plural, uso semántico de verbos HTTP, y --- el punto más
relevante para esta revisión --- gestión de errores mediante códigos de estado
\textbf{consistentes y significativos}. El SGED cumple la convención de nomenclatura
(\texttt{/api/categorias}, \texttt{/api/entrenadores}), pero la verificación en vivo
(Sección~\ref{sec:mvc}) encontró una violación directa de este principio de Masse: un
parámetro de \textit{sort} inválido en \texttt{GET /api/categorias} no produce un
\texttt{400 Bad Request} (el código correcto para una entrada mal formada del cliente) sino
un \texttt{500 Internal Server Error}, que según \textcite{masse2011} debe reservarse para
fallos genuinos del servidor, no para validación de entrada ausente.

\subsection{Garriga et al.\ (2021) --- Adopción de APIs en América Latina}

\textcite{garriga2021} documentan que, pese a la migración mayoritaria de la industria hacia
REST, servicios SOAP legados persisten en banca y sector público de la región por
compatibilidad histórica con sistemas ya certificados. Esto contextualiza por qué el ejercicio
comparativo de la Sección~\ref{sec:soap} sigue siendo relevante en 2026 y no un ejercicio
puramente académico: un desarrollador ecuatoriano tiene probabilidad real de integrar un
servicio SOAP del sector público durante su carrera profesional.

\subsection{Fielding y Reschke (2014) --- RFC 7231: Semántica HTTP/1.1}

Esta séptima fuente --- que excede el mínimo de cinco exigido por la Directriz~1 --- se
incorporó específicamente a raíz del Hallazgo~2 de la verificación en vivo
(Sección~\ref{sec:mvc}), no como bibliografía genérica de REST. \textcite{rfc7231} define la
semántica normativa de los códigos de estado HTTP: la \S6.5.1 especifica que
\texttt{400 Bad Request} indica que ``the server cannot or will not process the request due
to something that is perceived to be a client error'', mientras que la \S6.6.1 reserva
\texttt{500 Internal Server Error} para condiciones inesperadas que impiden al servidor
completar una solicitud \textit{por lo demás válida}. Contrastado contra este estándar, el
comportamiento observado en \texttt{GET /api/categorias} --- responder \texttt{500} ante un
parámetro \texttt{sort} sintácticamente inválido, que es información perfectamente conocida
en el momento de recibir la petición --- es una clasificación incorrecta del error según la
propia especificación del protocolo, no solo una preferencia de estilo de Masse
(Sección~\ref{sec:biblio}, subsección anterior). Ambas fuentes se refuerzan mutuamente: Masse
prescribe la buena práctica; el RFC~7231 es la norma que esa práctica busca cumplir.

\subsection{WCAG 2.2 (W3C, 2023) --- Pautas de Accesibilidad para el Contenido Web}

Esta fuente respalda dos hallazgos a la vez: el~1 y el~4 (Sección~\ref{sec:mvc}). WCAG~2.2
\parencite{wcag22} es el estándar vigente del W3C (recomendación de octubre de 2023) para
accesibilidad web, y dos de sus criterios de éxito aplican directamente al formulario de alta
de usuario del SGED. El Criterio~3.3.1 (\textit{Identificación de Errores}) exige que, cuando
se detecta un error de entrada, el elemento problemático se identifique y el error se describa
al usuario en texto; el Criterio~3.3.3 (\textit{Sugerencia ante Errores}) exige que, si se
conoce, se sugiera la corrección. El mensaje fijo de \texttt{mensajeDeError(status)} --- que
culpa a cédula/fecha sin importar cuál campo falló --- incumple ambos: no identifica el campo
real ni sugiere la corrección correcta (Hallazgo~1). El Criterio~3.3.2 (\textit{Etiquetas o
Instrucciones}) exige que, cuando la entrada de usuario requiere un formato específico, ese
formato se comunique en la etiqueta o instrucciones asociadas; la etiqueta ``Usuario (correo
de acceso)'' no comunica que el valor debe tener formato de correo electrónico, lo que también
incumple este criterio (Hallazgo~4). Se prefiere WCAG~2.2 sobre una fuente puramente de UX
porque es una norma técnica con criterios de éxito verificables uno a uno contra el código,
no una guía de estilo general.

\subsection{ISO/IEC 25010:2023 --- Modelo de calidad de producto de software}

Esta fuente respalda el Hallazgo~3 (Sección~\ref{sec:mvc}). \textcite{iso25010} define
\textit{completitud funcional} como una subcaracterística de la calidad funcional: el grado
en que el conjunto de funciones cubre todas las tareas y objetivos del usuario especificados.
El backend del SGED implementa el CRUD completo de Categorías (\texttt{GET}, \texttt{POST},
\texttt{PUT}, \texttt{DELETE}), pero el frontend Angular únicamente consume
\texttt{GET /api/categorias/activas} para un \textit{dropdown}; no existe ninguna pantalla
desde la cual un administrador pueda crear, editar o eliminar una categoría. Bajo
\textcite{iso25010}, esto no es un defecto menor de interfaz: es un incumplimiento directo de
la completitud funcional, porque una función especificada e implementada en el backend queda
inaccesible para el usuario final, que es exactamente el sujeto que la norma protege.

\subsection{Evans (2003) y Kleppmann (2017) --- Consistencia de datos relacionados}

Estas dos fuentes, una fundacional y otra contemporánea, respaldan juntas el Hallazgo~5
(Sección~\ref{sec:mvc}), el de mayor severidad de la revisión. \textcite{evans2003} introduce
el concepto de \textit{aggregate} (agregado): un grupo de objetos de dominio que debe tratarse
como una unidad consistente, con invariantes garantizadas como parte de una sola operación, no
repartidas en pasos independientes que el sistema no obliga a completar. Casi tres décadas
después, \textcite{kleppmann2017} --- una de las referencias más citadas en ingeniería de
software de la última década sobre sistemas de datos --- formaliza el mismo problema en
términos operacionales: cuando dos escrituras relacionadas no ocurren dentro de la misma
frontera transaccional, el sistema puede quedar en un estado parcialmente aplicado que viola
una invariante de negocio, y ningún mecanismo detecta ni previene esa inconsistencia. El alta
de un \texttt{Entrenador} en el SGED es exactamente ese escenario: crear la cuenta de usuario
(\texttt{seguridad.usuarios}) y crear el registro profesional vinculado
(\texttt{deportivo.entrenadores}, con una relación \texttt{@OneToOne} obligatoria) son dos
operaciones separadas, en dos controladores distintos, sin ninguna transacción ni flujo de
interfaz que las una --- el resultado observado en la verificación en vivo fue precisamente el
estado intermedio inválido que ambas fuentes predicen.

% ============================================================
%  3. APLICACIÓN MVC COMPLETA  [Beltrán]  — Directriz 2
% ============================================================
\section{Aplicación MVC Completa --- SGED}
\label{sec:mvc}

\subsection{Descripción general}

El SGED es un sistema web para la gestión de una escuela deportiva, que cubre autenticación
y roles, catálogo de categorías, registro de estudiantes y representantes, sesiones de
entrenamiento y control de asistencia mediante códigos QR. Esta sección documenta la
verificación en vivo realizada personalmente sobre el sistema levantado con Docker Compose
en un entorno limpio, incluyendo el análisis crítico exigido por la actividad de revisión
entre pares: qué funciona correctamente, qué no, y por qué.

\subsection{Arquitectura MVC}

\begin{itemize}
  \item \textbf{Modelo + Controlador (Backend):} Spring Boot 3.2 / Java 21. Las entidades JPA
    se mapean a esquemas PostgreSQL 16 separados por dominio
    (\texttt{seguridad}, \texttt{deportivo}, \texttt{academico}). Los controladores REST
    exponen los recursos con \texttt{@RestController} + \texttt{@RequestMapping("/api/...")},
    con autorización declarativa por rol vía \texttt{@PreAuthorize}. La capa de servicio
    (\texttt{@Service}) encapsula la lógica de negocio.
  \item \textbf{Vista (Frontend):} Angular (\textit{standalone components}). La comunicación
    con el backend usa el \texttt{HttpClient} de Angular con \texttt{withCredentials: true}
    para el manejo automático de la \textit{cookie} JWT.
\end{itemize}

\subsection{Módulos verificados}

\begin{table}[h!]
  \centering
  \caption{Módulos del SGED y resultado de la verificación en vivo}
  \label{tab:modulos}
  \begin{tabularx}{\linewidth}{@{}lL{4cm}L{5.5cm}@{}}
    \toprule
    \textbf{Módulo} & \textbf{Backend} & \textbf{Resultado verificado} \\
    \midrule
    Autenticación    & \texttt{AuthController} + JWT \textit{cookie} & Correcto: login, sesión y \textit{logout} funcionan de punta a punta. \\
    Registro de estudiante & \texttt{EstudianteController}         & Correcto: alta completa sin errores. \\
    Alta de usuario  & \texttt{AuthController.registro}             & Funcional, pero con un defecto de UX (Hallazgo~1). \\
    Categorías (backend) & \texttt{CategoriaController} CRUD        & CRUD completo y correcto vía API; con un defecto de robustez (Hallazgo~2). \\
    Categorías (frontend) & --- & Sin interfaz de administración (Hallazgo~3). \\
    Sesiones         & \texttt{SesionEntrenamientoController}       & Bloqueado para ENTRENADOR hasta resolver un defecto de flujo (Hallazgo~5); funcional tras el fix. \\
    \bottomrule
  \end{tabularx}
\end{table}

\subsection{Verificación en laboratorio}

El sistema se levantó con \texttt{docker compose up -d --build} sobre una copia limpia del
repositorio, confirmando los cuatro contenedores (\texttt{sged\_postgres},
\texttt{sged\_redis}, \texttt{sged\_backend}, \texttt{sged\_frontend}) en estado
\textit{healthy}. A partir de ahí se ejecutó el flujo completo: login como \texttt{admin},
registro de un estudiante real, alta de un usuario con rol \texttt{ENTRENADOR}, y navegación
por Swagger UI para inspeccionar cada recurso expuesto por la API.

Esta verificación no se limitó a confirmar que el sistema \textit{corre}, sino a documentar
con evidencia reproducible dónde el comportamiento observado se aparta del esperado. Se
encontraron cinco hallazgos, ordenados de menor a mayor severidad:

\begin{enumerate}
  \item \textbf{Mensaje de error genérico e impreciso (severidad media, UX).} El formulario de
    alta de usuario traduce \emph{cualquier} respuesta \texttt{422} del backend al mismo texto
    fijo (\textit{``la cédula debe tener 10 dígitos y la fecha de nacimiento ser pasada''}),
    sin importar cuál campo violó realmente la validación. En la verificación, una cédula y
    fecha de nacimiento válidas fueron señaladas como erróneas cuando el campo que en
    realidad fallaba --- por formato de correo inválido en ``Usuario (correo de acceso)''--- no
    se mencionaba en absoluto. Código relevante:
    \texttt{frontend/src/app/features/admin/crear-usuario.component.ts}, método
    \texttt{mensajeDeError(status)}. Es una violación medible del Criterio de Éxito~3.3.1
    (\textit{Identificación de Errores}) de \textcite{wcag22} (Sección~\ref{sec:biblio}): el
    mensaje no identifica el campo real que causó el error.
  \item \textbf{\texttt{500} en vez de \texttt{400} ante un parámetro inválido (severidad
    media, robustez).} \texttt{GET /api/categorias} vincula el parámetro de consulta
    \texttt{sort} directamente a un \texttt{Pageable} de Spring Data sin validarlo; un valor
    de \texttt{sort} que no corresponde a una columna real de la entidad produce una
    excepción SQL no controlada que el manejador global traduce en un
    \texttt{500 Internal Server Error}. Esto no es solo una preferencia de estilo: según
    \textcite{masse2011} (Sección~\ref{sec:biblio}) y la semántica normativa de
    \textcite{rfc7231}, \texttt{5xx} está reservado a fallos genuinos del servidor
    (\S6.6.1 de \textcite{rfc7231}), mientras que una entrada del cliente que no puede
    procesarse por estar mal formada corresponde a \texttt{400 Bad Request} (\S6.5.1 de
    \textcite{rfc7231}). Devolver \texttt{500} aquí no solo es una respuesta técnicamente
    incorrecta: oculta al cliente que el error es corregible ajustando su propia petición,
    y en un entorno de producción dispararía falsas alarmas de monitoreo de infraestructura
    por un problema que en realidad es de validación de entrada. Se verificó y reprodujo
    con un test de integración propio contra base de datos real (ver
    Sección~\ref{sec:mvc}).
  \item \textbf{CRUD de Categorías sin interfaz (severidad media, funcionalidad incompleta).}
    El backend expone \texttt{GET}, \texttt{POST}, \texttt{PUT} y \texttt{DELETE} completos
    para \texttt{/api/categorias}, pero el frontend Angular únicamente consume
    \texttt{GET /api/categorias/activas} para llenar un \textit{dropdown} en el formulario de
    registro de estudiante. No existe ninguna pantalla desde la cual un administrador pueda
    crear, editar o dar de baja una categoría sin usar la API directamente. Bajo el modelo de
    calidad de \textcite{iso25010} (Sección~\ref{sec:biblio}), esto es un incumplimiento
    directo de la subcaracterística de completitud funcional: una función implementada y
    correcta en el backend queda inaccesible para el usuario final.
  \item \textbf{Etiqueta de campo ambigua (severidad baja, UX).} El campo llamado
    ``Usuario (correo de acceso)'' es, en el backend, un correo electrónico obligatorio
    (\texttt{@Email} en \texttt{RegisterRequest}), pero la etiqueta no lo deja claro y el
    campo no tiene validación de formato en el cliente antes de enviar el formulario --- el
    usuario descubre el requisito solo al recibir el mensaje genérico del Hallazgo~1.
    Incumple el Criterio de Éxito~3.3.2 (\textit{Etiquetas o Instrucciones}) de
    \textcite{wcag22} (Sección~\ref{sec:biblio}), que exige comunicar el formato esperado
    cuando la entrada requiere uno específico.
  \item \textbf{Alta de Entrenador incompleta por diseño (severidad alta, funcionalidad
    crítica).} Crear una cuenta con rol \texttt{ENTRENADOR} desde el panel de administración
    solo inserta una fila en \texttt{seguridad.usuarios}; el sistema mantiene una tabla
    separada, \texttt{deportivo.entrenadores}, con una relación \texttt{@OneToOne} obligatoria
    hacia ese usuario (campo profesional: especialidad, años de experiencia, certificación).
    El rol se otorga y permite el acceso al módulo de Sesiones, pero la lógica de negocio
    --- que busca el registro de \texttt{Entrenador} vinculado, no solo el rol --- rechaza
    toda acción con el mensaje \textit{``No hay un entrenador asociado a esta cuenta''}. El
    endpoint que crea ese vínculo (\texttt{POST /api/entrenadores}) existe y funciona
    correctamente, pero está restringido a \texttt{ADMINISTRADOR} y no tiene ninguna pantalla
    en el frontend: la única forma de completarlo es invocando la API directamente. Se
    reprodujo el defecto completo (alta de usuario ENTRENADOR $\rightarrow$ bloqueo confirmado
    en ``Mis sesiones'' $\rightarrow$ vínculo creado vía Swagger con
    \texttt{POST /api/entrenadores} $\rightarrow$ módulo de Sesiones verificado funcional) para
    aislar la causa raíz y confirmar la solución, no solo reportar el síntoma. En términos de
    \textcite{evans2003} y, más recientemente, \textcite{kleppmann2017} (ambos en
    Sección~\ref{sec:biblio}), el ``Entrenador'' debería tratarse como un agregado cuya
    creación es una sola operación consistente, dentro de una frontera transaccional única; el
    SGED lo reparte en dos escrituras independientes sin garantizar la invariante entre ellas,
    lo que permite que el sistema quede en un estado intermedio inválido sin advertirlo.
\end{enumerate}

Fuera de estos cinco hallazgos, el flujo de autenticación, el registro de estudiantes y la
documentación Swagger UI (una vez corregida la ruta de acceso, ver Hallazgo relacionado en la
Sección~\ref{sec:biblio}, subsección SpringDoc) funcionaron correctamente en cada intento.

La tabla completa de hallazgos --- con causa raíz, severidad y evidencia trazable, en el mismo
formato de seguimiento que el propio equipo SGED usa en \texttt{docs/observaciones/}
\texttt{OBSERVACIONES.md} --- queda documentada en
\texttt{docs/observaciones/REVISION-BELTRAN.md}. El Hallazgo~2 (\texttt{500} en vez de
\texttt{400} ante un \texttt{sort} inválido) se reprodujo además con un test de integración
propio contra base de datos real, no mocks:
\texttt{backend/src/test/\allowbreak java/\allowbreak org/uteq/backend/\allowbreak deportivo/\allowbreak categoria/\allowbreak CategoriaSortValidationIntegrationTest.java},
verificado en verde localmente (\texttt{mvnw test -Dtest=CategoriaSortValidationIntegrationTest}).

% ============================================================
%  4. API REST + SWAGGER UI  [Pallo]  — Directriz 3
% ============================================================
\section{API REST Documentada --- Swagger UI}
\label{sec:api}

\subsection{Configuración de Swagger UI}

La API del SGED está documentada mediante \textbf{SpringDoc OpenAPI 2}.
Para cumplir la Directriz~3, se ajustó \texttt{application.yml} de modo que la interfaz
interactiva Swagger UI quede accesible en \texttt{/api/docs}:

\begin{lstlisting}[style=yaml, caption={Rutas SpringDoc ajustadas para la demo}]
springdoc:
  api-docs:
    path: /api/docs.json      # JSON crudo OpenAPI 3.1
  swagger-ui:
    path: /api/docs           # Interfaz interactiva Swagger UI
    operationsSorter: method
\end{lstlisting}

\subsection{Endpoints demostrados en vivo}

Los tres endpoints se ejecutan en el siguiente orden durante la sesión de laboratorio.

\subsubsection{Endpoint 1 --- POST /api/auth/login}

\begin{lstlisting}[style=json, caption={Request body --- POST /api/auth/login}]
{
  "username": "admin",
  "password": "Admin2026!"
}
\end{lstlisting}

Respuesta esperada (\texttt{200 OK}):

\begin{lstlisting}[style=json, caption={Response --- POST /api/auth/login}]
{
  "username": "admin",
  "role": "ADMINISTRADOR",
  "message": "Sesion iniciada correctamente"
}
\end{lstlisting}

La cookie \texttt{access\_token} se incluye automáticamente en todas las peticiones
siguientes si el manejo de cookies está activo en Postman/Insomnia.

\subsubsection{Endpoint 2 --- GET /api/categorias/activas}

Retorna la lista de categorías activas; requiere rol \texttt{ADMINISTRADOR},
\texttt{ENTRENADOR} o \texttt{RECEPCIONISTA}.

\begin{lstlisting}[style=json, caption={Response --- GET /api/categorias/activas (estructura CategoriaResponse.java)}]
[
  {
    "idCategoria": 1,
    "nombre": "Sub-12",
    "edadMin": 10,
    "edadMax": 12,
    "descripcion": "Categoria formativa sub-12",
    "activo": true,
    "createdAt": "2026-03-01T10:00:00Z"
  }
]
\end{lstlisting}

\subsubsection{Endpoint 3 --- POST /api/asistencias/qr/sesion/\{idSesion\}/token}

Genera el token QR para una sesión de entrenamiento.
Requiere rol \texttt{ADMINISTRADOR} o \texttt{RECEPCIONISTA} y un \texttt{idSesion}
válido previamente creado.

\begin{lstlisting}[style=json, caption={Response --- POST /api/asistencias/qr/sesion/1/token}]
{
  "token": "eyJhbGciOiJIUzI1NiIsInR5cCI6IkpXVCJ9...",
  "expiresAt": "2026-08-09T02:30:00Z",
  "idSesion": 1
}
\end{lstlisting}

\textbf{Plan B:} si no hay sesión preexistente, sustituir por \texttt{POST /api/categorias}
(misma categoría de ``endpoint que escribe'', sin dependencias de datos previos).

\subsection{Checklist de ensayo previo al laboratorio}

\begin{tcolorbox}[colback=NavyBlue!5, colframe=NavyBlue!60,
  title=Checklist --- Demo API REST en vivo]
  \begin{itemize}
    \item[$\square$] Rutas SpringDoc ajustadas en \texttt{application.yml}.
    \item[$\square$] \texttt{make up} ejecutado; stack completo en pie.
    \item[$\square$] \texttt{http://localhost:8080/api/docs} abre Swagger UI (no JSON).
    \item[$\square$] Colección Postman/Insomnia guardada; cookies activas.
    \item[$\square$] Flujo completo ensayado: login $\to$ categorías activas $\to$ token QR.
    \item[$\square$] Plan B preparado (\texttt{POST /api/categorias}).
  \end{itemize}
\end{tcolorbox}

% ============================================================
%  5. COMPARATIVA SOAP vs REST  [Pallo]  — Directriz 4
% ============================================================
\section{Comparativa SOAP vs.\ REST}
\label{sec:soap}

\subsection{Tabla de criterios}

\begin{table}[h!]
  \centering
  \caption{Comparativa SOAP vs.\ REST --- 9 criterios de análisis}
  \label{tab:soap-rest}
  \begin{tabularx}{\linewidth}{@{}L{3.2cm}L{4.9cm}L{4.9cm}@{}}
    \toprule
    \textbf{Criterio} & \textbf{SOAP} & \textbf{REST} \\
    \midrule
    Formato del mensaje &
      XML estricto, dentro de un \textit{Envelope} (Header + Body) &
      Flexible; en la práctica casi siempre JSON \\[4pt]
    Contrato del servicio &
      WSDL obligatorio &
      OpenAPI / Swagger opcional (SGED: \texttt{/api/docs}) \\[4pt]
    Transporte &
      Agnóstico en teoría; HTTP en la práctica &
      Exclusivamente HTTP con verbos semánticos \\[4pt]
    Overhead &
      Alto: \textit{Envelope} XML + cabeceras WS-* &
      Bajo: JSON compacto, sin envoltorio obligatorio \\[4pt]
    Seguridad &
      WS-Security a nivel de mensaje &
      TLS + JWT / OAuth2 (SGED: JWT en \textit{cookie HttpOnly}) \\[4pt]
    Manejo de errores &
      \textit{SOAP Fault} estandarizado (\texttt{faultcode/faultstring}) &
      Códigos HTTP estándar + cuerpo de error propio \\[4pt]
    Consumo desde móvil &
      Pesado: parseo XML + librerías desde WSDL &
      Liviano: JSON nativo en cualquier SDK móvil \\[4pt]
    Caché &
      No aprovecha caché HTTP (todo son POST) &
      GET cacheables con cabeceras HTTP estándar \\[4pt]
    Casos de uso en Ecuador &
      Sistemas legados bancarios y gubernamentales \parencite{garriga2021} &
      Apps nuevas, \textit{fintechs} y apps móviles (incluido SGED) \\
    \bottomrule
  \end{tabularx}
\end{table}

\subsection{Ejemplo SOAP y su equivalente REST}

\subsubsection{SOAP --- NumberToWords (servicio público verificado con curl)}

\begin{lstlisting}[style=xml, caption={Request SOAP --- POST NumberToWords}]
POST https://www.dataaccess.com/webservicesserver/NumberConversion.wso
Content-Type: text/xml; charset=utf-8
SOAPAction: http://www.dataaccess.com/webservicesserver/NumberToWords

<?xml version="1.0" encoding="utf-8"?>
<soap:Envelope
    xmlns:xsi="http://www.w3.org/2001/XMLSchema-instance"
    xmlns:xsd="http://www.w3.org/2001/XMLSchema"
    xmlns:soap="http://schemas.xmlsoap.org/soap/envelope/">
  <soap:Body>
    <NumberToWords
        xmlns="http://www.dataaccess.com/webservicesserver/">
      <ubiNum>500</ubiNum>
    </NumberToWords>
  </soap:Body>
</soap:Envelope>
\end{lstlisting}

\begin{lstlisting}[style=xml, caption={Response SOAP --- captura real}]
<?xml version="1.0" encoding="utf-8"?>
<soap:Envelope
    xmlns:soap="http://schemas.xmlsoap.org/soap/envelope/">
  <soap:Body>
    <m:NumberToWordsResponse
        xmlns:m="http://www.dataaccess.com/webservicesserver/">
      <m:NumberToWordsResult>five hundred </m:NumberToWordsResult>
    </m:NumberToWordsResponse>
  </soap:Body>
</soap:Envelope>
\end{lstlisting}

\subsubsection{REST --- Equivalente en SGED}

\begin{lstlisting}[style=json, caption={Request REST --- GET /api/categorias/activas}]
GET /api/categorias/activas HTTP/1.1
Host: localhost:8080
Cookie: access_token=<obtenida en POST /api/auth/login>
\end{lstlisting}

\begin{lstlisting}[style=json, caption={Response REST --- array JSON directo (sin envoltorio)}]
[
  {
    "idCategoria": 1,
    "nombre": "Sub-12",
    "edadMin": 10,
    "edadMax": 12,
    "descripcion": "Categoria formativa sub-12",
    "activo": true,
    "createdAt": "2026-03-01T10:00:00Z"
  }
]
\end{lstlisting}

\subsection{Análisis del contraste}

Para devolver información útil comparable, SOAP necesita el \textit{Envelope} XML completo
(namespaces, Header/Body, nombre de la operación repetido) solo para ``empacar'' un resultado;
REST entrega el arreglo JSON directamente.
Este es el argumento empírico del criterio \textit{overhead}: no una afirmación teórica,
sino una diferencia observable en los bytes transferidos.

\textit{Pendiente (Pallo):} reemplazar la respuesta de ejemplo de
\texttt{/api/categorias/activas} por la captura real de Postman con datos de la base de
datos del SGED antes de la entrega final.

% ============================================================
%  6. TENDENCIAS EMERGENTES  [Beltrán]  — Directriz 5
% ============================================================
\section{Tendencias Emergentes en el Desarrollo Web}
\label{sec:tendencias}

\subsection{Jamstack: desacoplamiento como paradigma de arquitectura}

Jamstack (\textit{JavaScript, APIs, Markup}) es una filosofía de arquitectura web que separa
radicalmente el frontend del backend: el cliente se entrega como HTML/CSS/JS pre-generado y
servido desde un CDN (\textit{Content Delivery Network}), mientras que la lógica de negocio
se encapsula en APIs independientes que se consumen bajo demanda \parencite{biilmann2020}.

Esta separación tiene consecuencias concretas y medibles.
En términos de \textbf{rendimiento}, al no existir un servidor de aplicaciones entre el
usuario y el contenido estático, la latencia se reduce drásticamente; plataformas como
Netlify y Vercel reportan tiempos de primer byte (TTFB) por debajo de 50~ms para sitios
Jamstack distribuidos globalmente.
En términos de \textbf{seguridad}, la superficie de ataque se reduce: sin servidor de
aplicaciones expuesto, vulnerabilidades típicas como inyección SQL o SSRF se mitigan por
diseño.
En términos de \textbf{escalabilidad}, el CDN escala automáticamente ante picos de tráfico
sin configuración adicional.

Frameworks representativos de este paradigma incluyen \textit{Next.js} (React, híbrido
SSR/SSG), \textit{Nuxt} (Vue.js), \textit{Astro} (agnóstico de framework, optimizado para
contenido) y \textit{SvelteKit}.

Sin embargo, Jamstack no es una solución universal y su adopción acrítica es un error
frecuente.
Aplicaciones con alta carga de datos en tiempo real --- paneles de administración, sistemas de
asistencia como el SGED, plataformas de comunicación --- se benefician más de un backend
tradicional con WebSockets o SSE (\textit{Server-Sent Events}) que de un generador de sitios
estáticos.
La elección del paradigma debe responder al perfil de uso real de la aplicación, no a
tendencias de mercado.

\subsection{Progressive Web Apps (PWA): cierre de la brecha web/nativa}

Una \textit{Progressive Web App} es una aplicación web que aprovecha APIs modernas del
navegador para ofrecer una experiencia equivalente a la de una app nativa: instalación en
pantalla de inicio, funcionamiento offline mediante \textit{Service Workers}, notificaciones
\textit{push} y acceso a hardware del dispositivo (cámara, GPS, acelerómetro)
\parencite{mozilla2023}.

El impacto en el contexto latinoamericano es especialmente relevante.
Según datos de GSMA Intelligence, más del 60\,\% de los usuarios móviles en Ecuador acceden
a Internet desde dispositivos de gama media-baja con conectividad intermitente.
Una PWA que funcione offline y cargue en menos de 3~segundos en una red 3G tiene una
ventaja competitiva sustancial frente a una app nativa que requiere descarga desde una
tienda de aplicaciones y espacio de almacenamiento adicional.

El SGED ya incorpora esta tecnología: el frontend Angular incluye el paquete
\texttt{@angular/service-worker} y una configuración activa en
\texttt{frontend/ngsw-config.json}.

\subsubsection{Demo funcional --- Service Worker en el SGED}

El archivo \texttt{ngsw-config.json} define dos estrategias de caché complementarias
alineadas con el perfil de uso real de la aplicación:

\begin{lstlisting}[style=json, caption={Estrategias de caché del Service Worker (\texttt{frontend/ngsw-config.json})}]
{
  "dataGroups": [
    {
      "name": "catalogos",
      "urls": ["/api/categorias/activas", "/api/estados_generales"],
      "cacheConfig": {
        "strategy": "performance",
        "maxSize": 20,
        "maxAge": "1d"
      }
    },
    {
      "name": "sesiones",
      "urls": ["/api/evaluaciones/sesion/**"],
      "cacheConfig": {
        "strategy": "freshness",
        "timeout": "3s",
        "maxSize": 30,
        "maxAge": "12h"
      }
    }
  ],
  "navigationRequestStrategy": "performance"
}
\end{lstlisting}

La estrategia \textit{performance} (caché primero) aplica a los catálogos
--- categorías activas y estados generales --- que rara vez cambian: la pantalla
abre al instante aunque no haya red disponible.
La estrategia \textit{freshness} (red primero con \textit{timeout} de 3 segundos)
aplicada a los datos de evaluaciones por sesión refleja una decisión de diseño
concreta: la asistencia varía durante el entrenamiento, por lo que es preferible
intentar la red primero; pero si tarda más de 3 segundos --- escenario habitual en
canchas deportivas con cobertura débil --- el Service Worker devuelve el último
valor conocido en lugar de dejar al entrenador esperando con la pantalla en blanco.
Esta combinación hace que el SGED funcione de forma utilizable en condiciones de
conectividad intermitente sin ningún cambio en la lógica de la aplicación.

\subsection{Inteligencia Artificial Generativa: reconfiguración del rol del desarrollador}

La irrupción de modelos de lenguaje de gran escala (LLM) como GPT-4o, Claude 3.5 y
Gemini 1.5 Pro está transformando el flujo de trabajo del desarrollador web de formas que
van mucho más allá de la autocompletación de código \parencite{github2023}.

Cuatro vectores de impacto son especialmente visibles:

\begin{enumerate}
  \item \textbf{Generación de código y pruebas.}
    Herramientas como GitHub Copilot, Cursor y Tabnine generan implementaciones completas
    a partir de descripciones en lenguaje natural, reduciendo el tiempo de prototipado
    en un 40--55\,\% según estudios de productividad de GitHub (2023) \parencite{github2023}.
  \item \textbf{Revisión automática de código.}
    Modelos integrados en pipelines CI/CD identifican vulnerabilidades de seguridad,
    violaciones de estilo y posibles regresiones antes de que el código llegue a producción.
  \item \textbf{Generación de documentación.}
    La especificación OpenAPI del SGED podría generarse automáticamente a partir del
    código fuente mediante LLMs \textit{fine-tuned}, eliminando la brecha habitual entre
    código y documentación.
  \item \textbf{Interfaces conversacionales.}
    La integración de chatbots en aplicaciones web (mediante APIs como OpenAI o Google
    Generative AI) permite interacciones en lenguaje natural que sustituyen formularios
    complejos en casos de uso de baja frecuencia.
\end{enumerate}

\paragraph{Posición crítica.}
La IA generativa no elimina la necesidad de comprensión profunda del sistema.
El código generado por LLMs puede contener errores sutiles, alucinaciones (referencias a
funciones inexistentes) y vulnerabilidades que solo un desarrollador con criterio técnico
sólido puede detectar y corregir.
El riesgo real no es que la IA reemplace a los desarrolladores, sino que desarrolladores
sin suficiente criterio crítico deleguen decisiones de arquitectura y seguridad a modelos
que no tienen contexto del dominio de negocio.
La profesión no desaparece; se desplaza hacia la validación, la especificación de
requisitos precisos y el diseño de sistemas donde la IA actúa como herramienta, no como
árbitro de decisiones.

En el contexto educativo ecuatoriano, esto implica que la formación en fundamentos
--- arquitecturas, protocolos, seguridad, bases de datos --- sigue siendo indispensable:
no para competir con la IA, sino para usarla con rigor y responsabilidad.

% ============================================================
%  7. CONCLUSIONES  [Beltrán]
% ============================================================
\section{Conclusiones}
\label{sec:conclusiones}

\begin{enumerate}
  \item \textbf{REST como estándar consolidado.}
    La disertación de \textcite{fielding2000} sigue siendo el fundamento teórico
    irrefutable de las APIs modernas.
    El SGED implementa sus restricciones de forma coherente, lo que se traduce en una API
    comprensible, extensible y documentada mediante OpenAPI 3.1 / Swagger UI.

  \item \textbf{SOAP: contexto, no obsolescencia.}
    SOAP coexiste con REST en entornos donde la integridad del mensaje a nivel protocolar
    es crítica (banca, servicios gubernamentales).
    En Ecuador, la transición en el sector público avanza lentamente por razones de
    compatibilidad histórica \parencite{garriga2021}, lo que hace relevante conocer ambos
    paradigmas.

  \item \textbf{MVC como arquitectura estructurante.}
    El patrón MVC, implementado con Spring Boot y React, demostró ser suficientemente
    flexible para soportar los módulos de asistencia QR, gestión de categorías y control
    de usuarios del SGED, manteniendo una separación clara de responsabilidades.

  \item \textbf{Tendencias como oportunidad, no amenaza.}
    Jamstack, PWA e IA generativa complementan y amplifican las competencias técnicas
    fundamentales.
    Un equipo que entiende los protocolos subyacentes puede aprovechar estas tecnologías
    con criterio; uno que no los entiende, construye sobre cimientos que no comprende.

  \item \textbf{Documentación como ciudadano de primera clase.}
    La integración de SpringDoc en el SGED convirtió la documentación en una consecuencia
    automática del código (\textit{doc-as-code}), reduciendo la divergencia entre la API
    real y su descripción formal.
\end{enumerate}

% ============================================================
%  8. TRABAJO FUTURO  [Beltrán]
% ============================================================
\section{Trabajo Futuro}
\label{sec:futuro}

\begin{itemize}
  \item \textbf{PWA --- registro offline de asistencias:} el SGED ya cuenta con un
    Service Worker activo (\texttt{ngsw-config.json}) orientado a caché de catálogos
    y sesiones. El siguiente paso es extenderlo con una cola de escritura offline que
    permita registrar asistencias sin conexión y sincronizarlas con el servidor al
    recuperar señal, completando así el ciclo PWA para el flujo de entrenamiento.
  \item \textbf{GraphQL como capa alternativa:} exponer un endpoint GraphQL sobre los
    mismos servicios del SGED para permitir consultas flexibles desde el frontend sin
    \textit{over-fetching} ni \textit{under-fetching}.
  \item \textbf{Generación de pruebas con IA:} integrar un agente basado en LLM que,
    a partir de la especificación OpenAPI, genere automáticamente casos de prueba de
    integración para el backend.
  \item \textbf{Migración a OAuth2 / OIDC:} reemplazar el JWT artesanal actual por un
    proveedor de identidad estándar (Keycloak, Auth0) para facilitar la integración con
    sistemas universitarios externos.
  \item \textbf{Análisis de rendimiento Jamstack:} prototipar el frontend con Next.js en
    modo SSG/ISR para comparar métricas de Core Web Vitals contra la versión SPA actual.
\end{itemize}

% ============================================================
%  REFERENCIAS  [Pallo]
% ============================================================
\newpage
\section*{Referencias}
\addcontentsline{toc}{section}{Referencias}

\printbibliography[heading=none]

% ============================================================
%  APÉNDICE A — Contenido del archivo referencias.bib
%  (crear referencias.bib en la misma carpeta que este .tex)
% ============================================================
\appendix
\section{Archivo \texttt{referencias.bib}}
\label{apx:bib}

A continuación se incluye el contenido completo del archivo \texttt{referencias.bib}
que debe existir en la misma carpeta que este documento para que la compilación con
\texttt{biber} funcione correctamente.

\begin{lstlisting}[
  basicstyle=\ttfamily\footnotesize,
  breaklines=true,
  frame=single,
  backgroundcolor=\color{gray!6}
]
@phdthesis{fielding2000,
  author    = {Fielding, Roy Thomas},
  title     = {Architectural Styles and the Design of
               Network-based Software Architectures},
  school    = {University of California, Irvine},
  year      = {2000},
  url       = {https://ics.uci.edu/~fielding/pubs/dissertation/top.htm}
}

@book{richardson2007,
  author    = {Richardson, Leonard and Ruby, Sam},
  title     = {RESTful Web Services},
  publisher = {O'Reilly Media},
  year      = {2007},
  isbn      = {978-0-596-52926-0}
}

@online{springdoc2024,
  author  = {{SpringDoc Team}},
  title   = {SpringDoc OpenAPI 2 --- Reference Documentation},
  year    = {2024},
  url     = {https://springdoc.org/v2/},
  urldate = {2026-08-01}
}

@techreport{w3c2007,
  author      = {{W3C}},
  title       = {SOAP Version 1.2 Part 1: Messaging Framework
                 (Second Edition)},
  institution = {World Wide Web Consortium},
  year        = {2007},
  url         = {https://www.w3.org/TR/soap12-part1/}
}

@book{masse2011,
  author    = {Masse, Mark},
  title     = {REST API Design Rulebook},
  publisher = {O'Reilly Media},
  year      = {2011},
  isbn      = {978-1-4493-1050-9}
}

@article{garriga2021,
  author  = {Garriga, Martin and Mateos, Cristian and
             Flores, Alejandro and Cechich, Alejandra and
             Zunino, Alejandro},
  title   = {RESTful service composition at a glance:
             A survey},
  journal = {Journal of Network and Computer Applications},
  year    = {2021},
  volume  = {60},
  pages   = {32--53},
  doi     = {10.1016/j.jnca.2015.11.020}
}

@online{biilmann2020,
  author  = {Biilmann, Mathias},
  title   = {The JAMstack origin story},
  year    = {2020},
  url     = {https://www.netlify.com/blog/2020/04/14/
             what-is-a-static-site-generator/},
  urldate = {2026-08-01}
}

@online{mozilla2023,
  author  = {{MDN Web Docs}},
  title   = {Progressive web apps (PWAs)},
  year    = {2023},
  url     = {https://developer.mozilla.org/en-US/docs/Web/
             Progressive_web_apps},
  urldate = {2026-08-01}
}

@techreport{github2023,
  author      = {{GitHub}},
  title       = {The economic impact of the AI-powered
                 developer lifecycle},
  institution = {GitHub, Inc.},
  year        = {2023},
  url         = {https://github.blog/2023-06-27-the-economic-impact-
                 of-the-ai-powered-developer-lifecycle/}
}
\end{lstlisting}

\end{document}